\section{Introduction}

Improving retrieval does not automatically yield better tabular agents. An agent answering a question over a public-data lake routinely finds the right dataset in the first few search results, peeks the schema, fires a query that returns zero rows, retries with different filters, and exhausts its budget switching between candidates --- never having had a retrieval problem at all. The bottleneck is not finding sources; it is \emph{committing to them}.

Text-centric Agentic RAG~\cite{lewis2020rag,trivedi2023ircot,jiang2023flare,asai2024selfrag,yao2023react} was designed around a one-shot abstraction: a passage either is or is not relevant, and reading it costs roughly one tool call. Exploratory question answering over tabular data lakes~\cite{wang2026lakeqa} breaks that abstraction. A useful table requires schema inspection, file-handle resolution, code execution, parsing repair, and a deliberate stopping decision. Each candidate source represents a multi-turn commitment, not a single read.

Tabular exploratory QA therefore needs an \textbf{execution-grounded control plane}, not (only) better retrieval. The planner must know how to bound its commitment to each source, how to keep low-level wrangling artifacts (schemas, tracebacks, row dumps) out of its own reasoning context, and how to report partial progress when one source covers some but not all of the required outputs.

\paragraph{Failure mechanics in single-agent tabular QA.} Six failure mechanisms recur in traces of single-agent baselines. \emph{Multi-turn source commitment} is unsupported: standard agentic-RAG loops have no primitive for ``finish or release this source.'' \emph{Stopping-rule absence} produces re-search/re-peek/re-query oscillation. \emph{Context pollution} fills the planner's window with schemas and tracebacks that degrade reasoning on the original question. \emph{File-handle errors} waste turns on ``file not found'' retries when a dataset is referenced by logical ID instead of exact path. \emph{Source drift} occurs when peeking is cheap and the planner rationalizes detours into tangential datasets. \emph{Premature submit} produces partial answers because nothing forces the planner to enumerate what it has not yet found.

\paragraph{Tabular MOMO.} We introduce \textbf{MOMO} (Multi-Agent Orchestration for Multi-Source Observation), a contract-driven runtime control plane for exploratory tabular QA. MOMO has three roles. A \emph{planner} orchestrates without touching raw data tools. A \emph{search subagent} executes a \texttt{SearchContract} that bounds discovery under an explicit budget. An \emph{inspect subagent} executes an \texttt{InspectContract} scoped to declared source families and enumerated required outputs. Workers return compact payloads with a four-value status (\texttt{success}, \texttt{partial}, \texttt{failed}, \texttt{budget\_exhausted}), evidence pointers, and an enumerated \texttt{missing\_outputs} list. Schema clutter and execution failures stay inside the worker --- the \emph{context firewall}. Runtime guards enforce file-reference correctness and source-family boundaries.

The contributions of this workshop paper are systems contributions rather than architecture or accuracy claims:
\begin{enumerate}
\item \textbf{Problem insight.} Tabular exploratory QA fails because of execution discipline against messy sources, not retrieval. We catalogue six recurring failure mechanisms that motivate a control-plane response.
\item \textbf{System abstraction.} A \emph{source interaction} is represented as an explicit contract with required outputs, success criteria, evidence, missing outputs, retry hints, and budget. Contracts are the unit of bounded commitment.
\item \textbf{Runtime control plane.} A planner with a reduced tool surface plus bounded workers, runtime guards (file-reference, source-family), a compact-return protocol, and a context firewall isolates execution failures from planner reasoning.
\item \textbf{Workflow-level evaluation.} The harness records tool waste, context-token growth, source-switch churn, partial-success rates, and cost --- not only end-to-end accuracy --- and we present a failure-mode taxonomy mapped to control-plane mechanisms.
\end{enumerate}

The implementation lives under \texttt{sana\_evaluation/} for historical reasons; a rename is queued and does not affect any claim made here. Our companion work (SANA) isolates \emph{where} the bottleneck is via search/planning/execution ablations; MOMO is one control-plane intervention for the execution bottleneck SANA exposes.
